\documentclass[a4paper,man,natbib]{apa6}
\usepackage[english]{babel}
\usepackage[utf8x]{inputenc}
\usepackage{amsmath}
\usepackage{cancel}
\usepackage{graphicx}
\usepackage[colorinlistoftodos]{todonotes}

\title{VTOL tail-sitter V2}
\shorttitle{V2}
\author{NY}
\affiliation{Summer 2026}

\begin{document}
\maketitle

\section{Nomenclature}

\section{Assumptions}
\subsection{Constraints}
\subsubsection{Cost:\\}
The total component budget for this build is to fit under \$400.
\subsubsection{Dimensions:}
\begin{itemize}
    \item Wingspan of 54 inches/
    \item Maximum weight of 1.3kg, driven by hardware motor constrain
\item An aspect ratio of 14 is chosen to authentically represent Iteration 1's design
\end{itemize}
\subsubsection{Hardware:}
\begin{itemize}
    \item Main motor: KDE3510XF-475 (Max thrust of 22 N)
\end{itemize}

\section{Initial Sizing and aircraft specification calculations}

Assuming a reasonably efficient wing planform ($e \approx 0.8$), and an estimated initial zero-lift drag coefficient of $C_{d_0} = 0.03$, we have the folling basic aircraft metrics:
\begin{align*}
    b &= 1.3761 \text{ m}\\
    AR &= 14\\
    S &= 0.135 \text{ m}^2\\
    W &= 12.955 \text{ N}\\
    \rho &= 1.225 \textbf{ kg/m}^3\\
    k &= \frac{1}{\pi ARe} \approx 0.02842\\
\end{align*}

\subsection{Cruise}
To maximize endurance, the aircraft will fly at the minimum-drag velocity:
\begin{align*}
    v_{md} &= \sqrt{\frac{2n_zW}{\rho S C_{L_{md}}}}\\
    n_z &= 1\\
    C_{L_{md}} &= \sqrt{\frac{C_{d_0}}{k}} \approx 1.0274\\
    &\boxed{v_{md} = 12.188 \text{ m/s}}
\end{align*}

Thrust requirements, therefore follow as:
\begin{align*}
    T &= \frac{1}{2}\rho v^2 S C_D\\
    C_D &= C_{d_0} + kC_{L_{md}}^2 \approx 0.06\\
    &\boxed{T_{req} = 0.757 \text{ N}}
\end{align*}
\end{document}
